\chapter{Virtualizzazione e cloud computing}
\emergencystretch=2.5em

Nel capitolo precedente abbiamo visto le architetture che servono milioni di utenti: repliche, bilanciatori di carico,
microservizi. Tutte presuppongono di avere a disposizione \textbf{molte macchine}, da accendere, spegnere e spostare con
facilità. Questo capitolo racconta da dove vengono quelle macchine: la \textbf{virtualizzazione}, che trasforma un server
fisico in tanti server «virtuali», e il \textbf{cloud computing}, che vende quelle risorse come un servizio. È il punto di
vista del \textbf{deployment}: l'ultimo passo del System Design, in cui si decide su quale hardware e software gireranno
i sottosistemi (il diagramma di deployment del Capitolo 9).

% Stili delle figure di questo capitolo
\tikzset{
  strato/.style={rectangle, draw=#1!75!black, fill=#1!18, line width=0.7pt, font=\scriptsize\bfseries,
                 align=center, minimum height=0.55cm, inner sep=2pt},
  vmbox/.style={rectangle, rounded corners=2pt, draw=sworange!80!black, fill=sworange!10, line width=0.7pt},
  san/.style={cylinder, shape border rotate=90, aspect=0.25, draw=swblue!70!black, fill=swblue!14, line width=0.8pt,
              minimum width=1.3cm, minimum height=1.0cm, font=\scriptsize\bfseries, align=center},
  fr/.style={-{Stealth[length=2.2mm]}, line width=0.8pt, draw=swgray},
}
% \torre{(x,y)}{nome}{etichetta}: un server fisico disegnato come una torre
\newcommand{\torre}[3]{%
  \begin{scope}[shift={#1}]
    \draw[line width=0.6pt, draw=swgray!80!black, fill=swlight] (-0.3,-0.42) rectangle (0.3,0.42);
    \foreach \yy in {0.28,0.17,0.06} \draw[swgray, line width=0.5pt] (-0.2,\yy) -- (0.2,\yy);
    \fill[swgreen] (0.17,-0.3) circle (0.04);
    \node[font=\tiny, anchor=north, align=center] at (0,-0.44) {#3};
  \end{scope}
  \coordinate (#2) at #1;}

% ============================================================
\section{Perché la virtualizzazione}

\dfn{Virtualizzazione}{
  La \term{virtualizzazione} consiste nel creare una versione \textbf{virtuale} di una risorsa che normalmente è fornita
  fisicamente e appartiene a un sistema: un computer, un disco, una rete. È una tecnica per aumentare l'\textbf{efficacia}
  e l'\textbf{efficienza} dei data center complessi.
}

L'idea è vecchia: risale agli anni Sessanta, ai mainframe IBM che andavano divisi fra molti utenti. È tornata di moda una
quindicina di anni fa grazie a un cambiamento nei processori, che conviene raccontare dall'inizio.

\subsection{L'evoluzione delle CPU}

I microprocessori hanno sempre visto crescere rapidamente le prestazioni e calare i costi. Per decenni la strategia per
aumentare le prestazioni è stata una sola: aumentare la \textbf{frequenza} del processore, cioè la sua «velocità».

\Needspace{16\baselineskip}
\noindent\begin{minipage}[c]{0.42\textwidth}
\begin{itemize}
  \item 1978: Intel 8086, 4 MHz;
  \item 1985: Intel 386, 25 MHz (circa 6 volte tanto);
  \item 1995: Intel Pentium Pro, 200 MHz (circa 8 volte);
  \item 2005: Intel Pentium 4, 3,8 GHz (circa 19 volte);
  \item 2020: Intel i9-10980XE, 3,0 GHz.
\end{itemize}
Nel 2005 la corsa si \textbf{ferma}: quindici anni dopo un processore di fascia alta ha una frequenza addirittura più
bassa. Si erano raggiunti i limiti fisici del silicio (calore e consumi crescono troppo con la frequenza).
\end{minipage}\hfill
\begin{minipage}[c]{0.55\textwidth}
\centering
\begin{tikzpicture}
\begin{semilogyaxis}[width=7.6cm, height=5.0cm, xmin=1975, xmax=2023, ymin=2, ymax=10000,
    xtick={1978,1985,1995,2005,2020}, xticklabel style={/pgf/number format/1000 sep={}, font=\tiny},
    yticklabel style={font=\tiny}, ylabel={MHz}, ylabel style={font=\scriptsize},
    grid=major, grid style={swgray!20}]
  \addplot[line width=1.1pt, swblue, mark=*, mark options={fill=swblue}] coordinates
    {(1978,4) (1985,25) (1995,200) (2005,3800) (2020,3000)};
  \node[font=\tiny, anchor=south east, text=swred] at (axis cs:2005,3800) {stop nel 2005};
  \node[font=\tiny, anchor=north, text=swgray] at (axis cs:2020,2500) {3,0 GHz};
\end{semilogyaxis}
\end{tikzpicture}
\end{minipage}

La soluzione è stata \textbf{moltiplicare le unità di calcolo}: invece di un processore più veloce, più processori (i
\textbf{core}) sullo stesso chip. Nel 2020 l'Intel i9-10980XE a 3,0 GHz ha 18 unità di calcolo, l'AMD Ryzen Threadripper
3970X a 3,7 GHz ne ha 32.

\subsection{Il problema dei multi-core}

Tradizionalmente i programmi sono stati scritti per girare su un computer con \textbf{una sola CPU}, e la grande
maggioranza delle applicazioni è fatta di programmi \textbf{sequenziali}. Un programma sequenziale usa un core alla
volta, per quanti ne abbia la macchina. Come si sfruttano allora più unità in parallelo?

\begin{esempio}{Un programma sequenziale non usa i core in più}
Il programma somma i resti di tre miliardi di numeri due volte: la prima con uno \textit{stream} normale, la seconda con
uno \textit{stream} parallelo, che divide il lavoro fra tutti i core disponibili.

\begin{lstlisting}[style=java, name=Somma.java]
import java.util.stream.LongStream;

public class Somma {
  public static void main(String[] args) {
    int core = Runtime.getRuntime().availableProcessors();
    System.out.println("Unità di calcolo disponibili: " + core);

    long n = 3_000_000_000L;
    long t0 = System.nanoTime();
    long s1 = LongStream.rangeClosed(1, n).map(x -> x % 7).sum();            // un solo core
    long t1 = System.nanoTime();
    long s2 = LongStream.rangeClosed(1, n).parallel().map(x -> x % 7).sum(); // tutti i core
    long t2 = System.nanoTime();

    System.out.println("Stesso risultato: " + (s1 == s2));
    System.out.printf("Sequenziale: %d ms%n", (t1 - t0) / 1_000_000);
    System.out.printf("Parallelo:   %d ms%n", (t2 - t1) / 1_000_000);
  }
}
\end{lstlisting}

\begin{lstlisting}[style=shell, name=Terminale]
$ java Somma.java
Unità di calcolo disponibili: 4
Stesso risultato: true
Sequenziale: 5888 ms
Parallelo:   1301 ms
\end{lstlisting}
Su una macchina con 4 core la versione parallela è circa quattro volte più veloce. Ma la differenza l'ha fatta il
\textbf{programmatore}, che ha dovuto riscrivere il calcolo; la prima versione, quella che si scrive di solito, lascia tre
core su quattro a guardare.
\end{esempio}

Riscrivere ogni applicazione in forma parallela non è realistico. La virtualizzazione prende la strada opposta: invece di
dividere \textbf{un} programma su più core, fa girare sulla stessa macchina \textbf{tanti sistemi indipendenti}, ciascuno
con i propri programmi sequenziali. Per capire perché questo sia utile bisogna guardare com'era fatto un data center.

% ============================================================
\section{Il data center tradizionale}

Nel data center tradizionale ogni server offre \textbf{un solo servizio} o una sola applicazione (\textit{one workload,
one box}): un server per l'applicazione, uno per il dominio (l'autenticazione degli utenti), uno per il database, uno
per la posta, uno per il terminal server, uno per il fax. Un nuovo servizio significa un \textbf{nuovo server}, con il suo
storage.

\begin{center}
\begin{tikzpicture}
  \foreach \x/\nome/\cpu [count=\i] in {0/Application/12, 2.3/Dominio/8, 4.6/Database/15, 6.9/Posta/6,
                                       9.2/Terminal/10, 11.5/Fax/5} {
    \torre{(\x,0)}{s\i}{Server\\ \nome};
    \draw[swgray!70, line width=0.4pt] (\x+0.45,-0.42) rectangle (\x+0.65,0.42);
    \fill[swred!70] (\x+0.45,-0.42) rectangle (\x+0.65,-0.42+0.84*\cpu/100);
    \node[font=\tiny, anchor=west, text=swred] at (\x+0.66,-0.42+0.84*\cpu/100) {\cpu\%};
  }
  \node[note, anchor=west] at (-0.4,1.0) {ogni server consuma corrente, scalda, occupa spazio e va gestito fisicamente;
    la barra rossa è l'uso medio della CPU};
\end{tikzpicture}
\end{center}

\subsection{Le conseguenze}

\begin{itemize}
  \item \textbf{Risorse sottoutilizzate}. La maggior parte dei server usa in media fra il \textbf{5\% e il 15\%} della
        propria capacità di calcolo. Per i servizi critici serve inoltre la \textbf{ridondanza}: computer di riserva che
        normalmente non lavorano affatto.
  \item \textbf{Costi di installazione e gestione elevati}. Ogni server consuma molta elettricità, emette calore, occupa
        spazio e va gestito fisicamente; gli aggiornamenti su larga scala costano molto, e i backup vanno fatti
        macchina per macchina.
  \item \textbf{Servizi intrinsecamente interconnessi}. Il server di dominio, che autentica gli utenti, viene interrogato
        da tutti gli altri servizi. Che cosa succede se si guasta?
\end{itemize}

\Needspace{14\baselineskip}
\noindent\begin{minipage}[c]{0.5\textwidth}
\centering
\begin{tikzpicture}
  \torre{(0,1.6)}{a}{Application}
  \torre{(2.2,1.6)}{d}{Dominio}
  \torre{(4.4,1.6)}{b}{Database}
  \torre{(0,-0.4)}{p}{Posta}
  \torre{(2.2,-0.4)}{t}{Terminal}
  \torre{(4.4,-0.4)}{f}{Fax}
  \foreach \n in {a,b,p,t,f} \draw[fr, {Stealth[length=2mm]}-{Stealth[length=2mm]}, shorten <=0.45cm,
    shorten >=0.45cm] (\n) -- (d);
  \draw[swred, line width=1.6pt] ($(d)+(-0.4,-0.45)$) -- ($(d)+(0.4,0.45)$);
  \draw[swred, line width=1.6pt] ($(d)+(-0.4,0.45)$) -- ($(d)+(0.4,-0.45)$);
\end{tikzpicture}
\end{minipage}\hfill
\begin{minipage}[c]{0.47\textwidth}
Se un server si guasta, può bloccare i servizi degli altri server che comunicano con lui. Ne derivano:
\begin{itemize}
  \item l'interruzione delle operazioni legate a quel servizio;
  \item nei casi più gravi, l'interruzione di \textbf{tutte} le attività dell'azienda;
  \item costi in più (il fermo e il ripristino);
  \item fermi non sempre prevedibili.
\end{itemize}
\end{minipage}

% ============================================================
\section{Le macchine virtuali}

\subsection{L'idea}

Un server fisico è fatto di una serie di \textbf{risorse hardware}, di un \textbf{sistema operativo} e di uno o più
\textbf{programmi applicativi}. L'idea: se la CPU del server ha più unità di calcolo parallele, possiamo «simulare»
l'esistenza di \textbf{più computer}, indipendenti dall'hardware «reale» sottostante? Ciascuno avrebbe il suo sistema
operativo e i suoi programmi, e le unità di calcolo verrebbero finalmente usate.

\begin{center}
\begin{tikzpicture}
  % server fisico
  \node[font=\scriptsize\bfseries] at (1.5,3.15) {Server fisico};
  \node[strato=swgray, minimum width=3cm] (hw1) at (1.5,0) {Hardware (CPU, RAM, disco, rete)};
  \node[strato=swblue, minimum width=3cm] at (1.5,0.6) {Sistema operativo};
  \node[strato=sworange, minimum width=3cm, minimum height=1.5cm] at (1.5,1.77) {Applicazione};
  % server virtualizzato
  \node[font=\scriptsize\bfseries] at (8.2,3.15) {Server virtualizzato};
  \node[strato=swgray, minimum width=6.4cm] at (8.2,0) {Hardware (CPU, RAM, disco, rete)};
  \node[strato=swred, minimum width=6.4cm] at (8.2,0.6) {Hypervisor (livello di virtualizzazione)};
  \foreach \x/\so [count=\k] in {6.1/Linux, 8.2/Windows, 10.3/Linux} {
    \draw[vmbox] (\x-1.0,0.95) rectangle (\x+1.0,2.95);
    \node[font=\tiny\bfseries, text=sworange!70!black] at (\x,2.75) {VM \k};
    \node[strato=swblue, minimum width=1.8cm] at (\x,1.25) {\so};
    \node[strato=sworange, minimum width=0.55cm, font=\tiny] at (\x-0.6,1.95) {App};
    \node[strato=sworange, minimum width=0.55cm, font=\tiny] at (\x,1.95) {App};
    \node[strato=sworange, minimum width=0.55cm, font=\tiny] at (\x+0.6,1.95) {App};
  }
  \draw[fr, line width=1.2pt] (3.4,1.4) -- node[above, note]{virtualizzazione} (4.6,1.4);
\end{tikzpicture}
\end{center}

\dfn{Hypervisor}{
  L'\term{hypervisor} è un nuovo strato software che gira \textbf{prima} del sistema operativo, direttamente sopra
  l'hardware. Simula l'esistenza di più computer «virtuali» su un unico computer fisico: fornisce loro le risorse fisiche
  e le gestisce in modo concorrente a beneficio delle \term{macchine virtuali} (VM).
}

Su ogni computer virtualizzato si installa un sistema operativo \textbf{normale}, che non ha modo di sapere se sta girando
su un computer «vero» o virtuale. Dal suo punto di vista vede una CPU, della memoria, un disco e una scheda di rete: sono
risorse virtuali, ritagliate dall'hypervisor da quelle fisiche.

\subsection{Le proprietà di una macchina virtuale}

\begin{itemize}
  \item \textbf{Indipendenza dall'hardware}. L'hypervisor simula una configurazione di server \textbf{standard}, nascondendo
        l'hardware fisico reale. Di conseguenza le VM possono girare su server fisici diversi senza bisogno di essere
        riconfigurate.
  \item \textbf{Una VM è un file}. Una macchina virtuale è un file che contiene il sistema operativo, le applicazioni, i
        dati e la rappresentazione della memoria del server. Il file conserva lo stato del server virtuale e permette di
        copiarlo su un numero qualsiasi di server fisici. Esistono strumenti commerciali che trasformano il contenuto di un
        server fisico in una VM.
  \item \textbf{Isolamento}. Ogni VM è completamente isolata dalle altre VM che girano sullo stesso server fisico: se una
        subisce un attacco o va in crash, le altre non ne risentono.
  \item \textbf{Gestione centralizzata}. Tutte le VM di un data center si gestiscono da \textbf{un'unica finestra}: la
        gestione dell'intero data center è centralizzata, e l'amministrazione si semplifica molto.
\end{itemize}

\begin{esempio}{Creare e gestire una VM dalla riga di comando}
Con VirtualBox, un hypervisor gratuito che si può installare sul proprio portatile, una VM si crea e si configura con il
comando \code{VBoxManage}. Le risorse che le diamo sono numeri in un file di configurazione, non pezzi di hardware.

\begin{lstlisting}[style=shell, name=crea-vm.sh]
# crea e registra una nuova macchina virtuale
VBoxManage createvm --name web01 --ostype Ubuntu_64 --register
# le assegna 2 CPU virtuali e 4 GB di RAM, presi da quelli del server fisico
VBoxManage modifyvm web01 --cpus 2 --memory 4096
# crea il disco virtuale: un file che puo' crescere fino a 20 GB
DISCO="$HOME/VirtualBox VMs/web01/web01.vdi"
VBoxManage createmedium disk --filename "$DISCO" --size 20480
VBoxManage storagectl web01 --name SATA --add sata
VBoxManage storageattach web01 --storagectl SATA --port 0 --type hdd --medium "$DISCO"
# la avvia senza finestra, come si fa su un server
VBoxManage startvm web01 --type headless
\end{lstlisting}

Che una VM sia «un file» si vede guardando la sua cartella:
\begin{lstlisting}[style=shell, name=Terminale]
$ ls -lh ~/"VirtualBox VMs/web01/"
-rw------- 1 anna anna 6.1G web01.vdi     # il disco: occupa solo lo spazio usato davvero
-rw------- 1 anna anna 9.4K web01.vbox    # la configurazione (CPU, RAM, schede di rete), in XML
\end{lstlisting}
Copiare questi due file su un altro computer con VirtualBox significa \textbf{spostare il server}.
\end{esempio}

% ============================================================
\section{La virtualizzazione in esercizio}

\subsection{Un server o un'infrastruttura}

Nella virtualizzazione di un \textbf{singolo server} le macchine virtuali risiedono nello storage del server fisico: dal
basso verso l'alto ci sono il server fisico, l'hypervisor, le VM e, dentro ciascuna, il sistema operativo (Microsoft,
Linux\ldots) e le applicazioni (gestionali, office, database\ldots). Se però il server si guasta, si fermano tutte le sue
VM insieme: abbiamo concentrato in una macchina sola i rischi di sei.

\dfn{Virtualizzazione enterprise ad alta disponibilità}{
  Un data center enterprise deve avere almeno \textbf{due server fisici} e un'unità disco condivisa, la
  \term{SAN} (\textit{Storage Area Network}), collegati fra loro ad altissima velocità (fibra ottica o iSCSI). Nella
  virtualizzazione enterprise le macchine virtuali risiedono nello storage \textbf{SAN}, non nei server.
}

\Needspace{16\baselineskip}
\begin{center}
\begin{tikzpicture}
  \foreach \x/\n/\lab in {0/s1/Server 1, 4.6/s2/Server 2} {
    \draw[line width=0.8pt, draw=swgray!80!black, fill=swlight] (\x-1.4,0) rectangle (\x+1.4,1.6);
    \node[font=\scriptsize\bfseries, anchor=north] at (\x,-0.02) {\lab};
    \node[strato=swred, minimum width=2.6cm, font=\tiny\bfseries] at (\x,0.3) {hypervisor};
    \coordinate (\n) at (\x,0);
  }
  \foreach \x in {-0.9,0,0.9} \node[vmbox, minimum width=0.7cm, minimum height=0.75cm, font=\tiny] at (\x,1.05) {VM};
  \foreach \x in {3.7,4.6,5.5} \node[vmbox, minimum width=0.7cm, minimum height=0.75cm, font=\tiny] at (\x,1.05) {VM};
  \node[san] (san) at (2.3,-2.0) {SAN};
  \draw[line width=1.4pt, swteal] (s1) ++(0,-0.35) -- (san.north west);
  \draw[line width=1.4pt, swteal] (s2) ++(0,-0.35) -- (san.north east);
  \node[note, text=swteal] at (2.3,-0.85) {fibra o iSCSI};
  \node[note, anchor=west, text width=5.6cm] at (7.1,0.55) {I file delle VM stanno nella SAN, visibile a entrambi i
    server. Ogni server può quindi avviare \textbf{qualunque} VM: per eseguirla basta avere il suo file.};
\end{tikzpicture}
\end{center}

La migrazione da un data center tradizionale avviene \textbf{dal fisico al virtuale} (P2V, \textit{physical to virtual}):
i sei server della figura iniziale diventano sei VM sui due server dell'infrastruttura, con i loro dischi nella SAN.

\subsection{Come funziona l'alta disponibilità}

La virtualizzazione garantisce la \textbf{continuità operativa} grazie alla funzione di \term{alta disponibilità} (HA,
\textit{High Availability}): se un server si guasta, le sue macchine virtuali vengono \textbf{migrate immediatamente} sugli
altri server disponibili e riavviate. Il fermo si riduce al tempo necessario a riavviare le VM.

Le tecnologie più recenti offrono anche la \term{fault tolerance} (FT), che garantisce la \textbf{continuità} del
funzionamento di una VM dopo il guasto di un server, eliminando del tutto il fermo: è attualmente la soluzione migliore
per la continuità operativa. È il fail-over del Capitolo 12, applicato alle macchine invece che ai processi.

\subsection{La migrazione a caldo}

\dfn{Live migration}{
  La migrazione «a caldo» (\term{live migration}) sposta una VM \textbf{in esecuzione} da un server fisico a un altro,
  senza spegnerla. Permette di fare manutenzione su un server o di bilanciare il carico fra i server mantenendo operativi
  i server virtuali.
}

\begin{esempio}{Una migrazione a caldo con libvirt}
Su Linux l'hypervisor KVM si gestisce con \code{virsh}, il client di \textit{libvirt}. Dovendo spegnere \code{server1}
per manutenzione, si spostano le sue VM su \code{server2} mentre lavorano:
\begin{lstlisting}[style=shell, name=Terminale]
$ virsh list
 Id   Name     State
-----------------------
 1    web01    running
 2    db01     running

$ virsh migrate --live --persistent web01 qemu+ssh://server2/system
$ virsh migrate --live --persistent db01  qemu+ssh://server2/system
\end{lstlisting}
Durante la migrazione la memoria della VM viene copiata sull'altro server mentre la VM continua a girare; alla fine la VM
si ferma per una frazione di secondo, giusto il tempo di copiare le ultime pagine modificate. Il disco non va copiato:
sta nella SAN, che entrambi i server vedono.
\end{esempio}

% ============================================================
\section{I vantaggi della virtualizzazione}

\subsection{Il consolidamento}

Con la virtualizzazione si parla di \term{consolidamento} del data center: sfruttando al massimo le risorse di calcolo
parallele dei processori moderni, si riduce in modo significativo il numero di server necessari a fornire i servizi IT di
un'organizzazione.

\begin{itemize}
  \item \textbf{Ottimizzazione dell'hardware}. Sistemi operativi diversi girano contemporaneamente su una sola macchina, e
        la CPU passa dal 5--15\% di un server dedicato a un uso molto più alto.
  \item \textbf{Ottimizzazione dello spazio su disco}. A una VM viene allocato solo lo spazio disco di cui ha
        strettamente bisogno; il resto del volume resta disponibile per nuove VM.
\end{itemize}

\begin{center}
\begin{tikzpicture}
\begin{axis}[width=12cm, height=3.6cm, xmin=0, xmax=48, ymin=0, ymax=100, ytick={0,50,100},
    yticklabel={\pgfmathprintnumber{\tick}\%}, yticklabel style={font=\tiny}, xtick=\empty,
    xlabel={tempo (una giornata)}, xlabel style={font=\scriptsize}, legend style={font=\tiny, at={(0.99,0.97)},
    anchor=north east, draw=none, fill=white, fill opacity=0.8, text opacity=1}, legend cell align=left]
  \addplot[swgray, line width=0.9pt, domain=0:48, samples=97] {9 + 5*sin(deg(x/3)) + 3*sin(deg(x*1.7))};
  \addlegendentry{server fisico dedicato}
  \addplot[swred, line width=0.9pt, domain=0:48, samples=97] {62 + 12*sin(deg(x/3)) + 6*sin(deg(x*1.3))};
  \addlegendentry{lo stesso server, con molte VM}
\end{axis}
\end{tikzpicture}
\end{center}

\subsection{Semplificazione e costi}

\begin{itemize}
  \item \textbf{Backup e disaster recovery}. L'intero server (sistema operativo, applicazioni, dati, dispositivi e stato)
        non è altro che un file: fare copie di sicurezza per il disaster recovery diventa facile.
  \item \textbf{Indipendenza dall'hardware}, come visto sopra.
  \item \textbf{Rilasci semplificati}. Si possono \textbf{clonare} le VM «a caldo» per fare test e aggiornamenti senza
        compromettere l'integrità dell'originale.
  \item \textbf{Riduzione dei costi IT}. Meno server significa risparmiare energia, condizionamento e spazio fisico del
        data center.
  \item \textbf{Riduzione dei costi del personale}. Con la gestione centralizzata, da un'unica interfaccia di
        amministrazione, il reparto IT è più produttivo.
\end{itemize}

\begin{esempio}{Provare un aggiornamento senza rischi}
Prima di aggiornare il sistema operativo del server web, si fotografa lo stato della VM e se ne fa un clone su cui
provare l'aggiornamento. Se qualcosa va storto, si torna indietro in pochi secondi.
\begin{lstlisting}[style=shell, name=aggiornamento.sh]
# una "fotografia" dello stato attuale (disco e configurazione)
VBoxManage snapshot web01 take prima-dell-aggiornamento
# un clone indipendente su cui provare l'aggiornamento
VBoxManage clonevm web01 --name web01-test --register
VBoxManage startvm web01-test --type headless
# ... l'aggiornamento sul clone rompe qualcosa? Si butta via il clone:
VBoxManage unregistervm web01-test --delete
# e se l'avessimo fatto sull'originale, si torna alla fotografia
VBoxManage snapshot web01 restore prima-dell-aggiornamento
# il backup per il disaster recovery e' un file da copiare altrove
VBoxManage export web01 --output web01-backup.ova
\end{lstlisting}
\end{esempio}

\begin{esempio}{Il data center di una pubblica amministrazione a Napoli}
\begin{itemize}
  \item Circa \textbf{800 macchine virtuali} in esecuzione in parallelo, con bilanciamento del carico;
  \item circa \textbf{30 server fisici};
  \item \textbf{2 SAN}, con circa 500 TB di spazio;
  \item la gestione di 800 macchine da \textbf{2 console} di amministrazione.
\end{itemize}
Il rapporto di consolidamento è di circa $800/30 \approx 27$ VM per server fisico. Senza virtualizzazione sarebbero serviti
centinaia di server, ciascuno usato al 5--15\%.
\end{esempio}

\Needspace{18\baselineskip}
\begin{infobox}{{Macchine virtuali e container}}
  Il Capitolo 12 ha mostrato un'applicazione a microservizi in cui ogni servizio è un \textbf{container} Docker. Un
  container non è una VM: non contiene un sistema operativo intero, ma solo l'applicazione con le sue librerie, e
  \textbf{condivide il kernel} del sistema che lo ospita. È più leggero (si avvia in un secondo, pesa decine di MB invece
  di GB) e meno isolato. Le due tecniche convivono: nel cloud i container girano quasi sempre dentro macchine virtuali.
\begin{lstlisting}[style=docker, name=Dockerfile]
# si parte da un'immagine con Java gia' installato
FROM eclipse-temurin:21-jre
WORKDIR /app
# si aggiunge solo l'applicazione: niente sistema operativo
COPY target/negozio.jar app.jar
EXPOSE 8080
CMD ["java", "-jar", "app.jar"]
\end{lstlisting}
\end{infobox}

% ============================================================
\section{Il cloud computing}

\dfn{Cloud computing}{
  Nel \term{cloud computing} le risorse IT sono fornite \textbf{come un servizio}: potenza di calcolo, storage, database,
  code di messaggi. Si usano quando servono e si pagano per quanto si usano.
}

\begin{itemize}
  \item I cloud sfruttano le \textbf{economie di scala} dell'hardware di largo consumo: storage economico, reti ad alta
        banda, processori multicore.
  \item I data center sono \textbf{distribuiti geograficamente}.
  \item Le offerte principali sono quelle di Microsoft, Amazon, Google e pochi altri.
\end{itemize}

\begin{center}
\begin{tikzpicture}
  \node[cloud, cloud puffs=14, cloud puff arc=110, aspect=2.1, draw=swblue!60, fill=swblue!5, line width=0.8pt,
        minimum width=9cm, minimum height=4.2cm] (c) at (0,0) {};
  \node[chip=sworange, minimum width=5.6cm] at (0,0.95) {Application: monitoraggio, contenuti, collaborazione, finanza};
  \node[chip=swteal, minimum width=5.6cm] at (0,0.15) {Platform: identità, runtime, code, database};
  \node[chip=swblue, minimum width=5.6cm] at (0,-0.65) {Infrastructure: calcolo, storage a blocchi, rete};
  \foreach \ang/\dev in {150/portatili, 30/desktop, 205/telefoni, 335/tablet, 90/server} {
    \node[font=\tiny, text=swgray] at ($(c.center)+(\ang:5.4cm and 2.6cm)$) {\dev};
  }
\end{tikzpicture}
\end{center}

La virtualizzazione è la tecnica che rende possibile il cloud: un fornitore con centinaia di migliaia di server fisici li
divide in milioni di VM e le affitta. A cambiare, per chi le usa, è il \textbf{modello economico}: non si comprano server,
si noleggia capacità.

% ============================================================
\section{Scalabilità ed elasticità}

\dfn{Scalabilità ed elasticità}{
  \begin{itemize}
    \item La \term{scalabilità} è la capacità di un sistema, di una rete o di un processo di gestire quantità di lavoro
          crescenti in modo armonioso, o di essere ingrandito facilmente.
    \item L'\term{elasticità} è la capacità di far \textbf{crescere o ridurre} dinamicamente le risorse
          dell'infrastruttura secondo il bisogno, per adattarsi ai cambiamenti del carico in modo \textbf{autonomo},
          massimizzando l'uso delle risorse.
  \end{itemize}
}

La scalabilità (Capitolo 12) dice che \textbf{si può} crescere; l'elasticità aggiunge che si cresce e si \textbf{decresce}
da soli, seguendo il carico. Le due proprietà si ottengono con il \textbf{provisioning dinamico} e con un progetto
\textbf{multi-tenant}, in cui più clienti (\textit{tenant}) condividono la stessa infrastruttura restando separati.

\subsection{Il provisioning dinamico}

\dfn{Provisioning dinamico}{
  Nel \term{provisioning dinamico} le istanze di calcolo vengono allocate o rilasciate (\textit{provisioned}) da una
  console di amministrazione o da un'applicazione client, dall'amministratore dei server o della rete o da qualunque
  altro utente abilitato, invece di essere acquistate e installate.
}

Nel modello tradizionale la capacità si decide \textbf{in anticipo}, e si sbaglia in uno dei due modi:

\Needspace{14\baselineskip}
\begin{center}
\begin{tikzpicture}
\pgfplotsset{prov/.style={width=5.4cm, height=3.6cm, xmin=0, xmax=12, ymin=0, ymax=10, xtick=\empty, ytick=\empty,
  xlabel={tempo}, ylabel={risorse}, xlabel style={font=\tiny}, ylabel style={font=\tiny},
  title style={font=\scriptsize\bfseries}}}
\begin{axis}[prov, title={Sotto-provisioning}]
  \addplot[name path=dmax, draw=none, domain=0:12, samples=100] {max(5 + 3.5*sin(deg(x*1.05)), 6)};
  \addplot[name path=cap, swred, dashed, line width=1pt] coordinates {(0,6) (12,6)};
  \addplot[swred!35] fill between[of=dmax and cap];
  \addplot[swblue, line width=1pt, domain=0:12, samples=100] {5 + 3.5*sin(deg(x*1.05))};
  \node[font=\tiny, text=swred] at (axis cs:3.0,9.3) {domanda persa};
\end{axis}
\begin{axis}[prov, title={Sovra-provisioning}, xshift=5.6cm]
  \addplot[name path=dom, swblue, line width=1pt, domain=0:12, samples=100] {5 + 3.5*sin(deg(x*1.05))};
  \addplot[name path=cap, swred, dashed, line width=1pt] coordinates {(0,9.2) (12,9.2)};
  \addplot[swgray!30] fill between[of=dom and cap];
  \node[font=\tiny, text=swgray!80!black] at (axis cs:6.4,7.5) {risorse inutilizzate};
\end{axis}
\begin{axis}[prov, title={Provisioning dinamico}, xshift=11.2cm]
  \addplot[swblue, line width=1pt, domain=0:12, samples=100] {5 + 3.5*sin(deg(x*1.05))};
  \addplot[swred, dashed, line width=1pt, const plot] coordinates
    {(0,6) (0.5,8) (1.0,9) (2.0,8) (2.5,6) (3.0,3) (3.8,2) (4.6,3) (5.2,6) (5.8,8.5) (7.0,9) (7.6,7) (8.2,4) (9,2)
     (10,3) (10.8,6) (11.4,8.5) (12,8.5)};
\end{axis}
\end{tikzpicture}
\end{center}

\begin{itemize}
  \item Chi \textbf{sottostima} l'uso del sistema compra troppo poco (\textit{under-provisioning}): nei picchi la domanda
        supera la capacità, e si perdono ricavi; se succede spesso, gli utenti se ne vanno e la domanda stessa cala
        (si perdono gli \textbf{utenti}).
  \item Chi lo \textbf{sovrastima} compra troppo (\textit{over-provisioning}): la capacità copre sempre i picchi, ma per
        la maggior parte del tempo le risorse restano inutilizzate (e pagate).
\end{itemize}

La soluzione è allocare le risorse \textbf{dinamicamente}, per seguire le variazioni stagionali della domanda, le
differenze fra settori diversi e i picchi improvvisi dovuti a eventi straordinari (un saldo, una diretta, un esame con
mille iscritti che si prenotano nello stesso minuto).

\begin{esempio}{Una regola di elasticità, in codice e in configurazione}
Un \textit{autoscaler} misura periodicamente l'uso medio della CPU delle istanze e ne calcola il numero adatto. La
regola di Kubernetes è semplice: se l'uso è il doppio dell'obiettivo, servono il doppio delle istanze.
\begin{lstlisting}[style=java, name=Autoscaler.java]
public class Autoscaler {
  private final int min, max;           // limiti decisi da chi paga
  private final double cpuObiettivo;    // per esempio 60%

  public Autoscaler(int min, int max, double cpuObiettivo) {
    this.min = min;
    this.max = max;
    this.cpuObiettivo = cpuObiettivo;
  }

  /** Quante istanze servono perche' l'uso medio della CPU torni all'obiettivo. */
  public int istanzeDesiderate(int istanzeAttuali, double cpuMedia) {
    int desiderate = (int) Math.ceil(istanzeAttuali * cpuMedia / cpuObiettivo);
    return Math.max(min, Math.min(max, desiderate));
  }
}
\end{lstlisting}
Con \code{new Autoscaler(2, 20, 60)}: di giorno 4 istanze al 90\% diventano $\lceil 4 \cdot 90/60 \rceil = 6$; di notte
6 istanze al 15\% diventerebbero $\lceil 6 \cdot 15/60 \rceil = 2$, il minimo. In un cloud non si scrive questo codice: lo
si \textbf{dichiara}, e il fornitore lo esegue.
\begin{lstlisting}[style=yaml, name=autoscaler.yaml]
apiVersion: autoscaling/v2
kind: HorizontalPodAutoscaler
metadata:
  name: negozio-web
spec:
  scaleTargetRef:            # che cosa scalare: il servizio web del negozio
    apiVersion: apps/v1
    kind: Deployment
    name: negozio-web
  minReplicas: 2             # mai meno di due, per la disponibilita'
  maxReplicas: 20            # mai piu' di venti, per il conto a fine mese
  metrics:
    - type: Resource
      resource:
        name: cpu
        target:
          type: Utilization
          averageUtilization: 60
\end{lstlisting}
\end{esempio}

% ============================================================
\section{I modelli di servizio}

Il cloud si offre a tre livelli, impilati uno sull'altro: l'infrastruttura (\textbf{IaaS}), la piattaforma
(\textbf{PaaS}) e il software (\textbf{SaaS}). Salendo, il fornitore si prende cura di una parte sempre più grande della
pila, e il cliente di una parte sempre più piccola.

\Needspace{16\baselineskip}
\begin{center}
\begin{tikzpicture}[x=2.55cm, y=0.44cm]
  \foreach \nome [count=\i] in {Applicazioni, Dati, Runtime, Middleware, Sistema operativo, Virtualizzazione, Server,
                                 Storage, Rete} {
    \node[font=\scriptsize, anchor=east] at (0.05,-\i) {\nome};
  }
  % colonne: in sede, IaaS, PaaS, SaaS; numero di strati gestiti dal cliente
  \foreach \col/\tit/\cli [count=\c] in {0/In sede/9, 1/IaaS/5, 2/PaaS/2, 3/SaaS/0} {
    \node[font=\scriptsize\bfseries] at (\col+0.6,0) {\tit};
    \foreach \i in {1,...,9} {
      \ifnum\i>\cli
        \fill[swblue!55] (\col+0.12,-\i+0.42) rectangle (\col+1.08,-\i-0.42);
      \else
        \fill[sworange!70] (\col+0.12,-\i+0.42) rectangle (\col+1.08,-\i-0.42);
      \fi
    }
  }
  \fill[sworange!70] (0.2,-10.6) rectangle (0.4,-10.2); \node[font=\scriptsize, anchor=west] at (0.42,-10.4) {gestito dal cliente};
  \fill[swblue!55] (1.9,-10.6) rectangle (2.1,-10.2); \node[font=\scriptsize, anchor=west] at (2.12,-10.4) {gestito dal fornitore};
\end{tikzpicture}
\end{center}

\subsection{Infrastructure as a Service (IaaS)}

\dfn{IaaS}{
  Al cliente viene data la possibilità di allocare \textbf{elaborazione, storage, reti} e altre risorse di calcolo
  fondamentali, su cui può installare ed eseguire software arbitrario, compresi sistemi operativi e applicazioni. Il
  cliente non gestisce né controlla l'infrastruttura cloud sottostante, ma ha il controllo dei \textbf{sistemi operativi},
  dello storage, delle applicazioni installate ed eventualmente un controllo limitato di alcuni componenti di rete.
}
Esempi: Amazon EC2, Eucalyptus, OpenNebula.

\begin{center}
\begin{tikzpicture}[bx/.style={line width=0.7pt, draw=#1!75!black, fill=#1!18}, tx/.style={font=\scriptsize\bfseries, align=center}]
  \draw[bx=sworange] (-4,1.0) rectangle node[tx]{Resource Management\\ Interface} (-0.1,1.75);
  \draw[bx=sworange] (0.1,1.0) rectangle node[tx]{System Monitoring\\ Interface} (4,1.75);
  \draw[bx=swred] (-4,0.35) rectangle node[tx]{Virtualization Layer} (4,0.85);
  \draw[bx=swgray] (-4,-0.55) rectangle node[tx]{Computing\\ System} (-1.4,0.2);
  \draw[bx=swgray] (-1.25,-0.55) rectangle node[tx]{Network\\ System} (1.25,0.2);
  \draw[bx=swgray] (1.4,-0.55) rectangle node[tx]{Storage\\ System} (4,0.2);
  \node[rotate=90, font=\tiny\bfseries, text=swblue] at (4.3,0.6) {Infrastructure};
\end{tikzpicture}
\end{center}

La tecnica che rende possibile l'IaaS è la \textbf{virtualizzazione}: un'astrazione delle risorse logiche rispetto alle
risorse fisiche sottostanti. Sposta il sistema operativo sopra l'hypervisor; più sistemi operativi condividono
l'hardware fisico e offrono servizi diversi, migliorando utilizzo, disponibilità, sicurezza e comodità. È esattamente la
figura del server virtualizzato vista all'inizio del capitolo: l'IaaS affitta le VM.

\begin{esempio}{Due VM in affitto}
Con la riga di comando di AWS si noleggiano due macchine virtuali EC2. Scegliamo noi l'immagine del sistema operativo e
la taglia; da lì in poi il sistema operativo è nostro: aggiornamenti, Java, sicurezza.
\begin{lstlisting}[style=shell, name=Terminale]
# due VM con 2 vCPU e 2 GB di RAM, partendo da un'immagine Ubuntu
$ aws ec2 run-instances --image-id ami-0abcdef1234567890 \
      --instance-type t3.small --count 2 --key-name chiave-ssh
# ci si collega come a un server qualsiasi, e si installa cio' che serve
$ ssh ubuntu@203.0.113.10 'sudo apt install -y openjdk-21-jre'
# quando non servono piu' si restituiscono, e si smette di pagare
$ aws ec2 terminate-instances --instance-ids i-0a1b2c3d4e5f60718 i-0f1e2d3c4b5a69788
\end{lstlisting}
\end{esempio}

\subsection{Platform as a Service (PaaS)}

\dfn{PaaS}{
  Al cliente viene data la possibilità di installare (\textit{deploy}) sull'infrastruttura cloud applicazioni create da
  lui o acquistate, scritte con i linguaggi e gli strumenti supportati dal fornitore. Il cliente non gestisce né
  controlla l'infrastruttura sottostante (rete, server, sistemi operativi, storage), ma ha il controllo delle
  \textbf{applicazioni} installate ed eventualmente della configurazione dell'ambiente che le ospita.
}
Esempi: Microsoft Azure, Google App Engine, Hadoop.

\begin{center}
\begin{tikzpicture}[bx/.style={line width=0.7pt, draw=#1!75!black, fill=#1!18}, tx/.style={font=\scriptsize\bfseries, align=center}]
  \draw[bx=sworange] (-4.5,1.05) rectangle (0.5,2.05);
  \node[tx] at (-2.0,1.83) {Programming IDE};
  \draw[bx=sworange, fill=white] (-4.4,1.15) rectangle node[font=\tiny\bfseries]{Programming APIs} (-2.05,1.6);
  \draw[bx=sworange, fill=white] (-1.95,1.15) rectangle node[font=\tiny\bfseries]{Development Tools} (0.4,1.6);
  \draw[bx=sworange] (0.7,1.05) rectangle (4.5,2.05);
  \node[tx] at (2.6,1.7) {System Control Interface};
  \node[font=\tiny] at (2.6,1.33) {policy, workflow};
  \draw[bx=swred] (-4.5,-0.3) rectangle (4.5,0.9);
  \node[tx, text=swred!60!black] at (0,0.7) {Runtime Environment};
  \foreach \x/\t in {-3.4/Scalability\\Elasticity, -1.15/Availability\\Reliability, 1.15/Performance\\Optimization,
                    3.4/Manageability\\Interoperability}
    \node[ellipse, draw=swred!70, fill=white, font=\tiny, align=center, inner sep=0.3pt, minimum width=2.05cm] at (\x,0.18) {\t};
  \draw[bx=swgray] (-4.5,-1.0) rectangle node[tx]{Computing} (-1.6,-0.45);
  \draw[bx=swgray] (-1.45,-1.0) rectangle node[tx]{Network} (1.45,-0.45);
  \draw[bx=swgray] (1.6,-1.0) rectangle node[tx]{Storage} (4.5,-0.45);
  \node[rotate=90, font=\tiny\bfseries, text=swteal] at (4.8,0.52) {Platform};
\end{tikzpicture}
\end{center}

La tecnica che rende possibile il PaaS è il progetto dell'\textbf{ambiente di esecuzione} (\textit{runtime environment}):
l'insieme dei servizi software disponibili, di solito realizzato da una collezione di librerie. Le sue proprietà tipiche
sono la gestibilità e l'interoperabilità, le prestazioni e l'ottimizzazione, la disponibilità e l'affidabilità, la
scalabilità e l'elasticità. Gli utenti sviluppano i propri servizi con un \textbf{IDE} che deve integrare tutte le
funzionalità dell'ambiente di esecuzione e offrire strumenti di sviluppo come profiler, debugger e ambienti di test. Le
API di programmazione cambiano da un fornitore all'altro, ma le operazioni comuni restano: calcolo, storage e
comunicazione.

\begin{esempio}{Consegnare il codice, non le macchine}
Con Google App Engine non si noleggiano VM: si consegna l'applicazione e un file che descrive l'ambiente che le serve.
Sistema operativo, Java, bilanciamento del carico e numero di istanze sono affare del fornitore.
\begin{lstlisting}[style=yaml, name=app.yaml]
runtime: java21                  # l'ambiente di esecuzione lo fornisce la piattaforma
instance_class: F2
automatic_scaling:               # l'elasticita' si chiede, non si costruisce
  min_instances: 1
  max_instances: 10
  target_cpu_utilization: 0.65
\end{lstlisting}
\begin{lstlisting}[style=shell, name=Terminale]
$ mvn package                    # si compila l'applicazione, come sempre
$ gcloud app deploy              # e la si consegna alla piattaforma
\end{lstlisting}
Rispetto all'esempio IaaS non c'è più nessun \code{ssh} né \code{apt install}: il confine fra ciò che gestiamo noi e ciò
che gestisce il fornitore si è alzato di due strati.
\end{esempio}

\subsection{Software as a Service (SaaS)}

\dfn{SaaS}{
  Al cliente viene data la possibilità di usare le \textbf{applicazioni del fornitore}, che girano su un'infrastruttura
  cloud. Le applicazioni sono accessibili da vari dispositivi attraverso un'interfaccia leggera, come un browser (per
  esempio la posta via web). Il cliente non gestisce né controlla l'infrastruttura né le singole funzionalità
  dell'applicazione, con l'eventuale eccezione di alcune impostazioni specifiche dell'utente.
}
Esempi: Google Apps (Gmail, Google Docs, Google Sites\ldots), Salesforce.com, EyeOS.

Le applicazioni tradizionali devono portare la propria interfaccia su una piattaforma \textbf{web}. Si trovano SaaS in
domini diversi:
\begin{itemize}
  \item applicazioni \textbf{generali}: suite per l'ufficio, multimedia, messaggistica istantanea;
  \item applicazioni \textbf{aziendali}: ERP, CRM, sistemi di compravendita;
  \item applicazioni \textbf{scientifiche}: simulazioni aerospaziali, simulazioni biochimiche;
  \item applicazioni per la \textbf{pubblica amministrazione}: sistemi sanitari nazionali, servizi di trasporto pubblico.
\end{itemize}

Per l'utente un SaaS è solo un sito. Per lo sviluppatore che deve integrarlo nel proprio sistema è spesso una
\textbf{API}: per esempio, il calendario di un utente si legge con una richiesta HTTP, senza installare nulla.
\Needspace{6\baselineskip}
\begin{lstlisting}[style=shell, name=Terminale]
$ curl -H "Authorization: Bearer $TOKEN" \
       https://www.googleapis.com/calendar/v3/users/me/calendarList
\end{lstlisting}
Il prossimo capitolo spiega come sono fatte queste API, e che cos'è quel token.

% ============================================================
\section{I modelli di deployment}

Come si dispiega un sistema cloud? Ci sono tre modelli principali: cloud \textbf{pubblico}, \textbf{privato} e
\textbf{ibrido}. Ciascuno può avere le caratteristiche viste finora; si differenziano soprattutto per la
\textbf{portata} e l'\textbf{accesso} ai servizi cloud, cioè per chi può usarli.

\subsection{Cloud pubblico}

\dfn{Cloud pubblico}{
  L'infrastruttura cloud è messa a disposizione del \textbf{pubblico} o di un grande gruppo industriale, ed è di
  proprietà di un'organizzazione che vende servizi cloud. È detto anche cloud esterno o \textit{multi-tenant}: un
  ambiente accessibile apertamente.
}
Caratteristiche: infrastruttura omogenea, politiche comuni, risorse condivise e multi-tenant, infrastruttura presa in
affitto, economie di scala.

\subsection{Cloud privato}

\dfn{Cloud privato}{
  L'infrastruttura cloud è gestita \textbf{solo per un'organizzazione}. Può essere amministrata dall'organizzazione stessa
  o da terzi, e trovarsi nella sede (\textit{on premise}) o altrove. Detto anche cloud interno, limita volutamente
  l'accesso alle risorse a chi appartiene all'organizzazione che lo possiede.
}
Caratteristiche: infrastruttura eterogenea, politiche personalizzate, risorse dedicate, infrastruttura interna, controllo
completo da un capo all'altro. Il data center della pubblica amministrazione di Napoli è, di fatto, un cloud privato.

\begin{center}
\renewcommand{\arraystretch}{1.25}
\begin{tabular}{lll}
\toprule
 & \textbf{Cloud pubblico} & \textbf{Cloud privato} \\
\midrule
Infrastruttura & omogenea & eterogenea \\
Politiche & definite in comune & personalizzate \\
Risorse & condivise, multi-tenant & dedicate \\
Costi & spesa operativa (si paga l'uso) & spesa in conto capitale (si compra) \\
Economia & grandi economie di scala & controllo completo \\
\bottomrule
\end{tabular}
\end{center}

\subsection{Cloud ibrido}

\dfn{Cloud ibrido}{
  L'infrastruttura è la \textbf{composizione} di due o più cloud (privati o pubblici) che restano entità distinte, ma sono
  legati da una tecnologia, standard o proprietaria, che permette di spostare dati e applicazioni fra l'uno e l'altro.
}

L'esempio tipico è il \term{cloud bursting}: un'azienda serve il carico normale con il proprio cloud privato e, nei
picchi, «trabocca» sul cloud pubblico, pagando la capacità in più solo quando serve. È il provisioning dinamico della
sezione precedente, applicato fra due cloud.

\begin{center}
\begin{tikzpicture}
  \node[cloud, cloud puffs=11, aspect=1.9, draw=swblue!60, fill=swblue!8, line width=0.8pt, minimum width=3.6cm,
        minimum height=1.6cm, font=\scriptsize\bfseries, text=swblue!70!black] (pub) at (5.2,1.1) {Cloud pubblico};
  \node[ellipse, draw=swgreen!60!black, fill=swgreen!10, line width=0.8pt, minimum width=4.2cm, minimum height=1.7cm,
        font=\scriptsize\bfseries, text=swgreen!50!black] (priv) at (0,0) {Cloud privato (azienda)};
  \draw[fr, line width=1.1pt, dashed, sworange] (priv.north east) to[bend left=20]
    node[above, note, text=sworange!80!black, align=center]{nei picchi: \textit{cloud bursting}} (pub.west);
  \node[note, anchor=north] at (0,-0.95) {carico normale, dati riservati};
  \node[note, anchor=north] at (5.2,0.25) {capacità in più, pagata a consumo};
\end{tikzpicture}
\end{center}

\subsection{L'ecosistema del cloud}

Mettendo insieme modelli di servizio e di deployment si ottiene la mappa del mercato: cloud pubblici, privati, ibridi e
\textbf{di comunità} (condivisi da organizzazioni con esigenze comuni), e a ogni livello fornitori diversi.

\begin{center}
\renewcommand{\arraystretch}{1.25}
\begin{tabular}{lP{10.5cm}}
\toprule
\textbf{Livello} & \textbf{Esempi di fornitori e prodotti} \\
\midrule
SaaS & Google Apps, Salesforce, Microsoft Dynamics CRM \\
PaaS & Microsoft Azure, Google App Engine, Force.com, IBM, Hadoop \\
IaaS & Amazon Web Services, Eucalyptus, GoGrid, Sun, AT\&T \\
\bottomrule
\end{tabular}
\end{center}

% ============================================================
\Needspace{24\baselineskip}
\gbox{In sintesi}{azure-gradient-6!80!black}{
\begin{itemize}
  \item Dal 2005 le CPU non diventano più veloci ma \textbf{più numerose}: i programmi sequenziali non sfruttano i core
        in più.
  \item Il data center tradizionale (un servizio per server) spreca risorse (5--15\% di CPU), costa molto e ha servizi
        interdipendenti: un guasto può fermare tutto.
  \item La \textbf{virtualizzazione} introduce l'\textbf{hypervisor}, che simula più computer su uno fisico. Una VM è
        indipendente dall'hardware, isolata, gestita centralmente, ed è \textbf{un file}.
  \item Con due server e una \textbf{SAN} si ottengono alta disponibilità (riavvio delle VM su un altro server), fault
        tolerance e \textbf{live migration}; con il consolidamento calano server, energia e personale.
  \item Il \textbf{cloud computing} vende le risorse IT come servizio. L'\textbf{elasticità} (crescere e decrescere da soli)
        si ottiene con il \textbf{provisioning dinamico}, contro il sotto- e il sovra-provisioning.
  \item Modelli di servizio: \textbf{IaaS} (si affittano VM: il sistema operativo è del cliente), \textbf{PaaS} (si
        consegna il codice), \textbf{SaaS} (si usa l'applicazione). Modelli di deployment: cloud \textbf{pubblico},
        \textbf{privato}, \textbf{ibrido} (con il cloud bursting).
\end{itemize}
}
